\chapter{Number Theory}

\begin{definition}[Divisibility, Primes, and Coprimality]
For $a, b\in\Z$, write $a\mid b$ if $b=ak$ for some $k\in\Z$. A prime is an integer $p>1$ whose only positive divisors are $1$ and $p$. Two integers $a, b$, not both zero, are coprime if $\gcd(a, b)=1$.
\end{definition}

\begin{definition}[Floor, Ceiling, and Fractional Part]
For $x\in\R$, the floor $\lfloor x\rfloor$ is the greatest integer at most $x$, the ceiling $\lceil x\rceil$ is the least integer at least $x$, and
\[
 \{x\}=x-\lfloor x\rfloor\in[0,1)
\]
is its fractional part. For every $m\in\Z$, $\lfloor x+m\rfloor=\lfloor x\rfloor+m$ and $\lceil x+m\rceil=\lceil x\rceil+m$.
\end{definition}
\begin{proposition}[Floor and Fractional-Part Identities]
For $x\in\R$ and every positive integer $m$,
\[
 \lfloor x\rfloor+\lfloor-x\rfloor=
 \begin{cases}-1,&x\notin\Z,\\0,&x\in\Z,\end{cases}
\qquad \lceil x\rceil=-\lfloor-x\rfloor,
\]
\[
 \sum_{j=0}^{m-1}\left\lfloor x+\frac jm\right\rfloor
 =\lfloor mx\rfloor.
\]
For all real $x,y$,
\[
 \lfloor x\rfloor+\lfloor y\rfloor\le\lfloor x+y\rfloor
 \le\lfloor x\rfloor+\lfloor y\rfloor+1.
\]
\end{proposition}

\begin{theorem}[Division Algorithm]
For $a\in\Z$ and $b\in\N$, there exist unique $q, r\in\Z$ such that
\[
 a=bq+r, \qquad 0\le r<b.
\]
\end{theorem}


\begin{proposition}[Euclidean Algorithm]
For integers $a, b$ with $b\ne0$,
\[
 \gcd(a, b)=\gcd(b, a\bmod|b|).
\]
Repeated Euclidean division produces the greatest common divisor as the last nonzero remainder.
\end{proposition}

\begin{theorem}[Fundamental Theorem of Arithmetic]
Every integer $n>1$ has a factorization
\[
 n=p_1^{\alpha_1}\cdots p_k^{\alpha_k},
\]
where the $p_i$ are distinct primes and $\alpha_i\in\N$. It is unique up to the order of the prime factors.
\end{theorem}

\begin{theorem}[Euclid's Lemma]
If a prime $p$ divides $ab$, then $p\mid a$ or $p\mid b$. More generally, if $\gcd(a, b)=1$ and $a\mid bc$, then $a\mid c$.
\end{theorem}

\begin{theorem}[B\'ezout's Identity]
For integers $a, b$, not both zero, there exist $x, y\in\Z$ such that
\[
 ax+by=\gcd(a, b).
\]
Moreover, the set of integer linear combinations $ax+by$ equals the set of multiples of $\gcd(a, b)$.
\end{theorem}

\begin{formula}[GCD and LCM from Prime Factorizations]
For positive integers $a, b$ and prime-power exponents $a=\prod_p p^{\alpha_p}$, $b=\prod_p p^{\beta_p}$ (allowing zero exponents),
\[
 \gcd(a, b)=\prod_p p^{\min(\alpha_p, \beta_p)},
 \qquad \lcmop(a, b)=\prod_p p^{\max(\alpha_p, \beta_p)}.
\]
Consequently,
\[
 \gcd(a, b)\lcmop(a, b)=ab.
\]
\end{formula}

\begin{theorem}[GCD of Differences of Powers]
For integers $a>1$ and $m, n\in\N$,
\[
 \gcd(a^m-1, a^n-1)=a^{\gcd(m, n)}-1.
\]
In particular, if $m\mid n$, then $a^m-1\mid a^n-1$.
\end{theorem}

\begin{theorem}[Euclid's Theorem]
There are infinitely many prime numbers.
\end{theorem}

\begin{theorem}[Bertrand's Postulate; Chebyshev's Theorem]
For every integer $n>1$, there is a prime $p$ satisfying
\[
 n<p<2n.
\]
\end{theorem}


\begin{definition}[Congruence and Residue Class]
For $n\in\N$,
\[
 a\equiv b\pmod n\quad\Longleftrightarrow\quad n\mid(a-b).
\]
The residue class of $a$ modulo $n$ is $[a]_n\in\Z/n\Z$. The group $(\Z/n\Z)^\times$ consists of the classes $[a]_n$ with $\gcd(a, n)=1$.
\end{definition}

\begin{proposition}[Algebra of Congruences]
If $a\equiv b\pmod n$ and $c\equiv d\pmod n$, then
\[
 a+c\equiv b+d\pmod n, \qquad
 ac\equiv bd\pmod n.
\]
If $ac\equiv bc\pmod n$ and $g=\gcd(c, n)$, then $a\equiv b\pmod{n/g}$; cancellation modulo $n$ is valid when $\gcd(c, n)=1$.
\end{proposition}

\begin{theorem}[Linear Congruence Criterion]
The congruence $ax\equiv b\pmod n$, with $n\ge1$, has an integer solution if and only if $d=\gcd(a, n)$ divides $b$. In that case there are exactly $d$ incongruent solutions modulo $n$.
\end{theorem}

\begin{theorem}[Chinese Remainder Theorem]
If $m_1, \dots, m_k\in\N$ are pairwise coprime, then for arbitrary integers $a_1, \dots, a_k$, the system
\[
 x\equiv a_i\pmod{m_i}\qquad(1\le i\le k)
\]
has a unique solution modulo $M=m_1\cdots m_k$.
\end{theorem}

\begin{proposition}[Chinese Remainder Theorem for General Moduli]
For positive moduli $m_1, \dots, m_k$, the congruences $x\equiv a_i\pmod{m_i}$ are simultaneously solvable if and only if
\[
 a_i\equiv a_j\pmod{\gcd(m_i, m_j)}
 \qquad\text{for all }i, j.
\]
When solvable, the solution is unique modulo $\lcmop(m_1, \dots, m_k)$.
\end{proposition}

\begin{formula}[Decimal Divisibility Tests]
For decimal digits $d_0, \dots, d_m$ and $N=\sum_{k=0}^m d_k10^k$,
\[
 N\equiv\sum_{k=0}^m d_k\pmod9, \qquad
 N\equiv\sum_{k=0}^m(-1)^k d_k\pmod{11}.
\]
The last $r$ digits determine $N\pmod{2^r}$ and $N\pmod{5^r}$.
\end{formula}

\begin{definition}[Positional Numeral Systems]
Let $b\ge2$ be an integer. Every nonnegative integer $N$ has a unique base-$b$ expansion
\[
 N=d_0+d_1b+\cdots+d_kb^k,
 \quad 0\le d_i<b,\quad d_k\ne0\ (N>0).
\]
The base-$b$ digit sum is $s_b(N)=\sum_{i=0}^k d_i$.
\end{definition}
\begin{formula}[Digit Sums and Repeated Digits]
For integers $b\ge2$ and $N\ge0$,
\[
 N\equiv s_b(N)\pmod{b-1},\qquad
 N\equiv\sum_{i=0}^k(-1)^id_i\pmod{b+1}.
\]
The base-$b$ integer consisting of $m$ digits all equal to $1$ is
\[
 1+b+\cdots+b^{m-1}=\frac{b^m-1}{b-1}.
\]
In base $10$, the last $k$ digits of $N$ are determined by $N\bmod 10^k$.
\end{formula}

\begin{theorem}[Rational and Eventually Periodic Expansions]
A real number has an eventually periodic base-$b$ expansion, for any fixed integer $b\ge2$, if and only if it is rational. A rational number in lowest terms has a terminating base-$b$ expansion if and only if every prime divisor of its denominator divides $b$.
\end{theorem}

\begin{formula}[Polynomial Congruences]
If $f\in\Z[x]$ and $a\equiv b\pmod n$, then $f(a)\equiv f(b)\pmod n$. In particular, the integral roots of $f(x)=0$ satisfy every reduction of that equation modulo $n$.
\end{formula}

\begin{theorem}[Wilson's Theorem]
An integer $n>1$ is prime if and only if
\[
 (n-1)!\equiv-1\pmod n.
\]
\end{theorem}


\begin{proposition}[Units and Inverses Modulo $n$]
An integer $a$ has a multiplicative inverse modulo $n$ if and only if $\gcd(a, n)=1$. If integers $u, v$ satisfy $au+nv=1$, then $u$ represents an inverse of $a$ modulo $n$.
\end{proposition}

\begin{definition}[Arithmetic and Multiplicative Functions]
An arithmetic function is a map $f\colon\N\to\C$. It is multiplicative if $f(1)=1$ and $f(mn)=f(m)f(n)$ whenever $\gcd(m, n)=1$. It is completely multiplicative if this identity holds for all $m, n\in\N$.
\end{definition}

\begin{definition}[Euler's Totient Function]
The Euler totient function is
\[
 \varphi(n)=|\{1\le k\le n\mid\gcd(k, n)=1\}|.
\]
It equals the order of the multiplicative group $(\Z/n\Z)^\times$.
\end{definition}

\begin{formula}[Euler's Totient Formula]
For $n=\prod_{i=1}^r p_i^{\alpha_i}$ with $\alpha_i\ge1$,
\[
 \varphi(n)=n\prod_{p\mid n}\left(1-\frac1p\right)
 =\prod_{i=1}^r p_i^{\alpha_i-1}(p_i-1).
\]
The empty product gives $\varphi(1)=1$. Also
\[
 \sum_{d\mid n}\varphi(d)=n.
\]
\end{formula}

\begin{formula}[Divisor Functions]
For $n=\prod_{i=1}^r p_i^{\alpha_i}$, the number $\tau(n)$ and sum $\sigma(n)$ of positive divisors satisfy
\[
 \tau(n)=\prod_{i=1}^r(\alpha_i+1), \qquad
 \sigma(n)=\prod_{i=1}^r\frac{p_i^{\alpha_i+1}-1}{p_i-1}.
\]
More generally, $\sigma_k(n)=\sum_{d\mid n}d^k=\prod_i(1+p_i^k+\cdots+p_i^{k\alpha_i})$ for $k\in\Nzero$.
\end{formula}

\begin{definition}[M\"obius Function]
The M\"obius function $\mu\colon\N\to\{-1, 0, 1\}$ is given by
\[
 \mu(n)=\begin{cases}
 1, &n=1, \\
 (-1)^r, &n\text{ is a product of $r$ distinct primes}, \\
 0, &n\text{ is divisible by a square greater than $1$}.
 \end{cases}
\]
\end{definition}

\begin{theorem}[M\"obius Inversion Formula]
For arithmetic functions $f, g$, if
\[
 g(n)=\sum_{d\mid n}f(d),
\]
then
\[
 f(n)=\sum_{d\mid n}\mu(d)g(n/d)
 =\sum_{d\mid n}\mu(n/d)g(d).
\]
In particular,
\[
 \sum_{d\mid n}\mu(d)=\begin{cases}1, &n=1, \\0, &n>1.\end{cases}
\]
\end{theorem}

\begin{formula}[GCD Sum and Totient Identity]
For every $n\in\N$,
\[
 \sum_{k=1}^n\gcd(k,n)
 =\sum_{d\mid n}d\,\varphi(n/d).
\]
More generally, grouping integers by their gcd with $n$ gives
\[
 \#\{1\le k\le n\mid\gcd(k,n)=d\}=\varphi(n/d)
 \qquad(d\mid n).
\]
\end{formula}

\begin{definition}[Dirichlet Convolution]
For arithmetic functions $f, g$, their Dirichlet convolution is
\[
 (f\ast g)(n)=\sum_{d\mid n}f(d)g(n/d).
\]
It is associative and commutative; the identity is $\varepsilon(1)=1$, $\varepsilon(n)=0$ for $n>1$. The constant function $\mathbf1(n)=1$ satisfies $\mu\ast\mathbf1=\varepsilon$.
\end{definition}


\begin{proposition}[Multiplicative Arithmetic Functions]
If $f, g$ are multiplicative arithmetic functions, then their Dirichlet convolution $f\ast g$ is multiplicative. Moreover, $\tau$, $\sigma_k$, $\varphi$, and $\mu$ are multiplicative.
\end{proposition}

\begin{theorem}[Fermat's Little Theorem]
If $p$ is prime and $p\nmid a$, then
\[
 a^{p-1}\equiv1\pmod p.
\]
Equivalently, $a^p\equiv a\pmod p$ for every integer $a$.
\end{theorem}

\begin{theorem}[Euler's Theorem]
For $n\in\N$ and $\gcd(a, n)=1$,
\[
 a^{\varphi(n)}\equiv1\pmod n.
\]
\end{theorem}

\begin{theorem}[Even Perfect Numbers; Euclid--Euler Theorem]
A positive even integer is perfect if and only if it has the form
\[
 2^{p-1}(2^p-1),
\]
where $p$ and $2^p-1$ are both prime.
\end{theorem}

\begin{definition}[Multiplicative Order]
For $\gcd(a, n)=1$, the multiplicative order $\ord_n(a)$ is the least positive integer $r$ such that $a^r\equiv1\pmod n$. It divides $\varphi(n)$, and
\[
 a^k\equiv1\pmod n
 \quad\Longleftrightarrow\quad\ord_n(a)\mid k
\]
for $k\in\Z$.
\end{definition}

\begin{proposition}[Repetend Length and Multiplicative Order]
Let $a/b$ be a reduced rational number with $b>1$ and $\gcd(b,B)=1$ for an integer base $B\ge2$. Its base-$B$ fractional expansion is purely periodic when $0\le a<b$, and its least positive period length is $\ord_b(B)$.
\end{proposition}

\begin{definition}[Primitive Root]
An integer $g$ is a primitive root modulo $n$ if $\gcd(g, n)=1$ and
\[
 \ord_n(g)=\varphi(n).
\]
Equivalently, $[g]_n$ generates $(\Z/n\Z)^\times$.
\end{definition}

\begin{theorem}[Primitive Root Theorem]
For $n>1$, primitive roots modulo $n$ exist exactly when
\[
 n\in\{2, 4, p^k, 2p^k\},
\]
where $p$ is an odd prime and $k\in\N$.
\end{theorem}

\begin{proposition}[Power Congruences in a Cyclic Group]
Suppose $(\Z/n\Z)^\times$ is cyclic of order $N=\varphi(n)$. For an integer $k\ge1$ and $\gcd(a, n)=1$, the congruence
\[
 x^k\equiv a\pmod n
\]
has solutions if and only if $a^{N/d}\equiv1\pmod n$, where $d=\gcd(k, N)$. When solvable, it has exactly $d$ solutions modulo $n$.
\end{proposition}

\begin{definition}[Carmichael Function]
For $n\ge1$, the Carmichael function $\lambda(n)$ is the exponent of $(\Z/n\Z)^\times$: the least positive integer $m$ such that $a^m\equiv1\pmod n$ for all $a$ coprime to $n$. For odd primes $p$,
\[
 \lambda(p^k)=p^{k-1}(p-1),
\]
while $\lambda(2)=1$, $\lambda(4)=2$, and $\lambda(2^k)=2^{k-2}$ for $k\ge3$. For pairwise coprime $a, b$, $\lambda(ab)=\lcmop(\lambda(a), \lambda(b))$.
\end{definition}


\begin{theorem}[Primitive Root Test]
Let $p$ be an odd prime and $g\in\Z$ with $p\nmid g$. Then $g$ is a primitive root modulo $p$ if and only if
\[
 g^{(p-1)/\ell}\not\equiv1\pmod p
\]
for every distinct prime divisor $\ell$ of $p-1$.
\end{theorem}

\begin{proposition}[The Group of Units Modulo a Power of Two]
For $k\ge3$,
\[
 (\Z/2^k\Z)^\times\cong C_2\times C_{2^{k-2}}.
\]
In particular, there is no primitive root modulo $2^k$ when $k\ge3$.
\end{proposition}

\begin{definition}[$p$-Adic Valuation]
For a prime $p$ and nonzero integer $n$, $v_p(n)$ is the largest $r\in\Nzero$ such that $p^r\mid n$. For nonzero integers $a, b$,
\[
 v_p(ab)=v_p(a)+v_p(b), \qquad
 v_p(a+b)\ge\min(v_p(a), v_p(b))
\]
when $a+b\ne0$. The second inequality is an equality if $v_p(a)\ne v_p(b)$.
\end{definition}

\begin{proposition}[Valuation of a Product and a Sum]
Let $p$ be prime and $x,y\in\Z\setminus\{0\}$. Then
\[
 v_p(xy)=v_p(x)+v_p(y),\qquad
 v_p(x+y)\ge\min(v_p(x),v_p(y))
\]
when $x+y\ne0$. If $v_p(x)\ne v_p(y)$, equality holds in the second inequality.
\end{proposition}

\begin{theorem}[Legendre's Formula]
For a prime $p$ and $n\in\N$,
\[
 v_p(n!)=\sum_{k\ge1}\left\lfloor\frac{n}{p^k}\right\rfloor
 =\frac{n-s_p(n)}{p-1},
\]
where $s_p(n)$ is the sum of the base-$p$ digits of $n$.
\end{theorem}

\begin{formula}[Legendre's Formula via Digit Sums]
For a prime $p$ and a nonnegative integer $n$,
\[
 v_p(n!)=\sum_{j\ge1}\left\lfloor\frac{n}{p^j}\right\rfloor
 =\frac{n-s_p(n)}{p-1}.
\]
Consequently, for nonnegative $r\le n$,
\[
 v_p\binom nr
 =\frac{s_p(r)+s_p(n-r)-s_p(n)}{p-1}.
\]
\end{formula}

\begin{theorem}[Lifting the Exponent Lemma; LTE]
Let $p$ be an odd prime, $n\in\N$, and $x, y\in\Z$ with $x\ne y$, $p\mid x-y$, and $p\nmid xy$. Then
\[
 v_p(x^n-y^n)=v_p(x-y)+v_p(n).
\]
If instead $p\mid x+y$, $p\nmid xy$, $n$ is odd, and $x^n+y^n\ne0$, then
\[
 v_p(x^n+y^n)=v_p(x+y)+v_p(n).
\]
For distinct odd integers $x, y$ and even $n$,
\[
 v_2(x^n-y^n)=v_2(x-y)+v_2(x+y)+v_2(n)-1.
\]
For odd $n$, $v_2(x^n-y^n)=v_2(x-y)$.
\end{theorem}

\begin{theorem}[Lucas's Theorem]
Let $p$ be prime and write $n=\sum_{i\ge0}n_ip^i$, $k=\sum_{i\ge0}k_ip^i$, with $0\le n_i, k_i<p$. For $0\le k\le n$,
\[
 \binom nk\equiv\prod_{i\ge0}\binom{n_i}{k_i}\pmod p,
\]
where $\binom{a}{b}=0$ if $b>a$.
\end{theorem}

\begin{theorem}[Kummer's Theorem]
For a prime $p$ and integers $0\le k\le n$, the valuation $v_p\!\left(\binom nk\right)$ equals the number of carries when $k$ and $n-k$ are added in base $p$.
\end{theorem}

\begin{formula}[Valuation of a Binomial Coefficient]
For a prime $p$ and $0\le k\le n$,
\[
 v_p\!\left(\binom nk\right)=
 \sum_{r\ge1}\left(
 \left\lfloor\frac n{p^r}\right\rfloor-
 \left\lfloor\frac k{p^r}\right\rfloor-
 \left\lfloor\frac{n-k}{p^r}\right\rfloor
 \right).
\]
\end{formula}


\begin{formula}[Powers of Two in Factorials]
For $n\in\N$,
\[
 v_2(n!)=n-s_2(n),
\]
where $s_2(n)$ is the number of ones in the binary expansion of $n$.
\end{formula}

\begin{proposition}[Valuations of Products]
For nonzero integers $a_1, \dots, a_m$,
\[
 v_p\!\left(\prod_{i=1}^m a_i\right)=\sum_{i=1}^m v_p(a_i).
\]
For nonzero integers $a, b$,
\[
 v_p(\gcd(a, b))=\min(v_p(a), v_p(b)),
\]
where $\gcd(a, b)$ is positive.
\end{proposition}

\begin{definition}[Quadratic Residue and Legendre Symbol]
For an odd prime $p$ and $a\in\Z$, $a$ is a quadratic residue modulo $p$ if $x^2\equiv a\pmod p$ for some integer $x$. The Legendre symbol is
\[
 \left(\frac ap\right)=
 \begin{cases}
 0, &p\mid a, \\
 1, &a\text{ is a nonzero quadratic residue modulo }p, \\
 -1, &\text{otherwise}.
 \end{cases}
\]
\end{definition}

\begin{theorem}[Euler's Criterion]
For an odd prime $p$ and any integer $a$,
\[
 a^{(p-1)/2}\equiv\left(\frac ap\right)\pmod p.
\]
If $p\nmid a$, there are exactly two solutions modulo $p$ to $x^2\equiv a\pmod p$ when the Legendre symbol is $1$, and none when it is $-1$.
\end{theorem}

\begin{proposition}[Multiplicativity of the Legendre Symbol]
For an odd prime $p$ and integers $a, b$,
\[
 \left(\frac{ab}{p}\right)
 =\left(\frac ap\right)\left(\frac bp\right).
\]
\end{proposition}


\begin{theorem}[Gauss's Lemma for Quadratic Residues]
Let $p$ be an odd prime, $a\in\Z$ with $p\nmid a$, and let $m$ be the number of integers $k\in\{1, \dots, (p-1)/2\}$ for which the least positive residue of $ak$ modulo $p$ exceeds $p/2$. Then
\[
 \left(\frac ap\right)=(-1)^m.
\]
\end{theorem}

\begin{theorem}[Law of Quadratic Reciprocity]
For distinct odd primes $p, q$,
\[
 \left(\frac pq\right)\left(\frac qp\right)
 =(-1)^{(p-1)(q-1)/4}.
\]
The supplementary laws are
\[
 \left(\frac{-1}{p}\right)=(-1)^{(p-1)/2}, \qquad
 \left(\frac2p\right)=(-1)^{(p^2-1)/8}.
\]
\end{theorem}

\begin{definition}[Jacobi Symbol]
For a positive odd integer $n=\prod_i p_i^{\alpha_i}$ and $a\in\Z$, define
\[
 \left(\frac an\right)
 =\prod_i\left(\frac a{p_i}\right)^{\alpha_i}.
\]
By convention $\left(\frac a1\right)=1$. For composite $n$, Jacobi symbol $1$ does not necessarily imply that $a$ is a quadratic residue modulo $n$.
\end{definition}

\begin{proposition}[Jacobi Symbol Identities]
For odd positive integers $m,n$ and integers $a,b$,
\[
 \left(\frac{ab}{n}\right)=\left(\frac an\right)\left(\frac bn\right),
\quad
 \left(\frac{a}{mn}\right)=\left(\frac am\right)\left(\frac an\right).
\]
If $\gcd(a,n)=1$, then $\left(\frac an\right)\in\{-1,1\}$; otherwise it is $0$. For odd $n$,
\[
 \left(\frac{-1}{n}\right)=(-1)^{(n-1)/2},\qquad
 \left(\frac{2}{n}\right)=(-1)^{(n^2-1)/8}.
\]
For coprime odd positive integers $m,n$,
\[
 \left(\frac mn\right)\left(\frac nm\right)
 =(-1)^{\frac{m-1}{2}\frac{n-1}{2}}.
\]
\end{proposition}

\begin{proposition}[Solutions of $x^2\equiv1$ Modulo Prime Powers]
For an odd prime $p$ and $k\ge1$, the solutions of $x^2\equiv1\pmod{p^k}$ are exactly $x\equiv\pm1\pmod{p^k}$. Modulo $2^k$, there is one solution for $k=1$, two for $k=2$, and four for $k\ge3$.
\end{proposition}


\begin{theorem}[Quadratic Residues of $-1$ and $2$]
For an odd prime $p$,
\[
 -1\text{ is a quadratic residue modulo }p
 \quad\Longleftrightarrow\quad p\equiv1\pmod4,
\]
\[
 2\text{ is a quadratic residue modulo }p
 \quad\Longleftrightarrow\quad p\equiv\pm1\pmod8.
\]
\end{theorem}

\begin{proposition}[Square and Cube Valuation Criteria]
An integer $N\ge1$ is a perfect square if and only if $v_p(N)$ is even for every prime $p$. More generally, $N$ is a positive perfect $k$th power if and only if $k\mid v_p(N)$ for every prime $p$. In particular, every integer square is congruent to $0$ or $1$ modulo $4$, and to $0$, $1$, or $4$ modulo $8$.
\end{proposition}
\begin{definition}[Diophantine Equation]
A Diophantine equation is an equation in which solutions are required to be integers (or, when explicitly specified, natural numbers or rational numbers). The domain restriction is part of the problem.
\end{definition}

\begin{theorem}[Linear Diophantine Equations]
Let $a, b, c\in\Z$, with $a, b$ not both zero. The equation $ax+by=c$ has integer solutions if and only if $d=\gcd(a, b)$ divides $c$. If $(x_0, y_0)$ is one solution, all solutions are
\[
 x=x_0+\frac bd t, \qquad y=y_0-\frac ad t,
 \qquad t\in\Z.
\]
\end{theorem}

\begin{theorem}[Primitive Pythagorean Triples]
All triples of positive integers $a, b, c$ satisfying $a^2+b^2=c^2$ and $\gcd(a, b, c)=1$ arise, after possibly interchanging $a, b$, as
\[
 a=m^2-n^2, \quad b=2mn, \quad c=m^2+n^2,
\]
where $m>n\ge1$, $\gcd(m, n)=1$, and $m, n$ have opposite parity.
\end{theorem}

\begin{theorem}[Frobenius Coin Theorem for Two Denominations]
If $a, b\ge2$ are coprime integers, the greatest nonnegative integer not expressible as $ax+by$ with $x, y\in\Nzero$ is
\[
 ab-a-b.
\]
The number of nonnegative integers not expressible in this form is $(a-1)(b-1)/2$.
\end{theorem}

\begin{theorem}[Fermat's Two-Square Theorem]
An odd prime $p$ can be written as a sum of two integer squares if and only if $p\equiv1\pmod4$. The prime $2=1^2+1^2$ is also such a sum.
\end{theorem}

\begin{theorem}[Sum of Two Squares Criterion]
A positive integer $n$ is a sum of two integer squares if and only if every prime $p\equiv3\pmod4$ appears to an even exponent in the prime factorization of $n$.
\end{theorem}

\begin{formula}[Brahmagupta--Fibonacci Identity]
For integers $a, b, c, d$,
\[
 (a^2+b^2)(c^2+d^2)=(ac-bd)^2+(ad+bc)^2.
\]
The companion expression $(ac+bd)^2+(ad-bc)^2$ has the same value.
\end{formula}

\begin{theorem}[Lagrange's Four-Square Theorem]
Every nonnegative integer is a sum of four integer squares.
\end{theorem}

\begin{theorem}[Legendre's Three-Square Theorem]
A nonnegative integer $n$ is a sum of three integer squares if and only if it is not of the form
\[
 n=4^a(8b+7),\qquad a,b\in\Nzero.
\]
\end{theorem}


\begin{proposition}[Solutions of Homogeneous Linear Diophantine Equations]
If $\gcd(a, b)=d$ and $a, b$ are not both zero, then every integer solution of $ax+by=0$ has the form
\[
 (x, y)=\left(\frac bd t, -\frac ad t\right),
 \qquad t\in\Z.
\]
\end{proposition}

\begin{theorem}[Zsigmondy's Theorem, Integer Form]
Let $a>b\ge1$ be coprime integers and $n>1$. Except when $n=2$ and $a+b$ is a power of $2$, or when $(a, b, n)=(2, 1, 6)$, the integer $a^n-b^n$ has a prime divisor which divides none of $a^k-b^k$ for $1\le k<n$.
\end{theorem}

\begin{theorem}[Dirichlet's Approximation Theorem]
For any irrational $\alpha\in\R$ and integer $N\ge1$, there exist integers $p$ and $q$ with $1\le q\le N$ such that
\[
 |q\alpha-p|<\frac1N.
\]
In particular, infinitely many reduced fractions $p/q$ satisfy $|\alpha-p/q|<1/q^2$.
\end{theorem}

\begin{definition}[Farey Sequences]
The Farey sequence $\mathcal{F}_N$ lists in increasing order all reduced fractions $a/b$ with $0\le a\le b\le N$. Consecutive terms $a/b<c/d$ satisfy $bc-ad=1$. Conversely, reduced fractions $a/b<c/d$ are neighbors in $\mathcal{F}_N$ if and only if $bc-ad=1$ and $b+d>N$.
\end{definition}

\begin{definition}[Simple Continued Fraction]
For integers $a_0\in\Z$, $a_i\in\N$ for $i\ge1$, the simple continued fraction
\[
 [a_0;a_1, \dots, a_n]
 =a_0+\cfrac1{a_1+\cfrac1{a_2+\dots+\cfrac1{a_n}}}
\]
has a reduced form $p_n/q_n$ with $q_n>0$, called its $n$th convergent.
\end{definition}

\begin{proposition}[Convergent Recurrences]
With initial values $p_{-2}=0$, $p_{-1}=1$, $q_{-2}=1$, $q_{-1}=0$, the convergents satisfy
\[
 p_n=a_np_{n-1}+p_{n-2}, \qquad
 q_n=a_nq_{n-1}+q_{n-2},
\]
and
\[
 p_nq_{n-1}-p_{n-1}q_n=(-1)^{n-1}.
\]
\end{proposition}

\begin{theorem}[Approximation by Continued Fractions]
If $\alpha\in\R\setminus\Q$ has convergents $p_n/q_n$, then
\[
 \left|\alpha-\frac{p_n}{q_n}\right|<\frac1{q_n^2}.
\]
\end{theorem}

\begin{theorem}[Legendre's Criterion for Convergents]
Let $\alpha$ be irrational and $p/q$ a reduced rational number with $q>0$. If
\[
 \left|\alpha-\frac pq\right|<\frac1{2q^2},
\]
then $p/q$ is a convergent of the simple continued fraction of $\alpha$.
\end{theorem}

\begin{theorem}[Best Approximation by Convergents]
Let $p_n/q_n$ be a convergent of an irrational number $\alpha$. For integers $p,q$ with $0<q<q_{n+1}$, one has
\[
 |q_n\alpha-p_n|\le |q\alpha-p|,
\]
with equality only when $(p,q)=(p_n,q_n)$ up to the possible sign convention for $q$. In particular, convergents are best approximations of the second kind.
\end{theorem}

\begin{definition}[Pell Equation]
For a fixed positive nonsquare integer $D$, Pell's equation is
\[
 x^2-Dy^2=1, \qquad x, y\in\Z.
\]
\end{definition}

\begin{theorem}[Solutions of Pell's Equation]
For every positive nonsquare integer $D$, there is a least solution $(x_1, y_1)$ in positive integers to $x^2-Dy^2=1$ with $x_1>1$. Every positive-integer solution $(x_n, y_n)$ is determined by
\[
 x_n+y_n\sqrt D=(x_1+y_1\sqrt D)^n,
 \qquad n\in\N.
\]
\end{theorem}

\begin{theorem}[Negative Pell Equation]
For a positive nonsquare integer $D$, the equation
\[
 x^2-Dy^2=-1
\]
has an integer solution if and only if the period length of the simple continued fraction of $\sqrt D$ is odd. If a least positive solution $(x_1,y_1)$ exists, all positive solutions are given by
\[
 x+y\sqrt D=(x_1+y_1\sqrt D)^{2k+1},\qquad k\in\Nzero.
\]
\end{theorem}


\begin{theorem}[Periodicity of Continued Fractions of Quadratic Irrationals]
A real irrational number has an eventually periodic simple continued fraction if and only if it is a quadratic irrational, that is, a root of an irreducible quadratic polynomial over $\Q$.
\end{theorem}

\begin{definition}[Cyclotomic Polynomial]
For $n\ge1$, the $n$th cyclotomic polynomial is
\[
 \Phi_n(x)=\prod_{\substack{1\le k\le n\\\gcd(k,n)=1}}
 \left(x-e^{2\pi i k/n}\right).
\]
It has integer coefficients, is monic of degree $\varphi(n)$, and satisfies
\[
 x^n-1=\prod_{d\mid n}\Phi_d(x).
\]
For a prime $p$, $\Phi_p(x)=1+x+\cdots+x^{p-1}$.
\end{definition}

\begin{definition}[Gaussian Integers and Norm]
The Gaussian integers form the ring $\Z[i]=\{a+bi\mid a, b\in\Z\}$. The norm is
\[
 N(a+bi)=a^2+b^2=(a+bi)(a-bi).
\]
It is multiplicative: $N(zw)=N(z)N(w)$. The units of $\Z[i]$ are $\pm1, \pm i$.
\end{definition}

\begin{theorem}[Unique Factorization in Gaussian Integers]
The ring $\Z[i]$ is a Euclidean domain with respect to its norm and hence a unique factorization domain. Every nonzero nonunit factors into Gaussian primes, uniquely up to order and multiplication by units.
\end{theorem}

\begin{theorem}[Finite Fields and Their Multiplicative Groups]
For each prime power $q$, there is up to isomorphism a unique field $\mathbb{F}_q$ with $q$ elements. Its nonzero elements form a cyclic group of order $q-1$, so every $a\in\mathbb{F}_q$ satisfies $a^q=a$.
\end{theorem}

\begin{theorem}[Hensel's Lemma, Simple-Root Form]
Let $p$ be prime and $f\in\Z[x]$. If an integer $a$ satisfies
\[
 f(a)\equiv0\pmod p,
 \qquad f'(a)\not\equiv0\pmod p,
\]
then for every $k\in\N$ there exists a unique residue class $x_k\pmod{p^k}$ with $x_k\equiv a\pmod p$ and $f(x_k)\equiv0\pmod{p^k}$.
\end{theorem}
